
\section{\method: Forward-Free Diffusion Language Model}

\label{sec:method}
To address these challenges in discrete diffusion, we formulate a forward-free diffusion model in which refinement is learned directly from sampler-induced intermediate drafts. 
We outline the proposed \method below and defer detailed derivations and proofs to Appendix~\ref{app:proofs}.

\subsection{Diffusion via Recursive Marginal Refinement}

\paragraph{Draft distribution initialization}
We begin by defining an initial draft distribution $p_{\theta,T}$. 
Since the goal is to approximate the target data distribution $p_0\rbr{\dv}$, a natural design principle is to make $p_{\theta,T}\rbr{\dv}$ as close as possible to $p_0\rbr{\dv}$, thereby reducing the number of refinement iterations required at inference time. 
This leads to the initial matching objective
\begin{equation}
    \min_{\theta}\KL\!\rbr{p_0(\dv)\,\|\,p_{\theta, T}(\dv)}.
    \label{eq:init_kl}
\end{equation}
If $p_{\theta,T}=p_0$, generation can be performed in a single step and no refinement is needed. 
In practice, finite model capacity, factorized decoding, and imperfect optimization introduce an approximation gap between $p_{\theta,T}$ and $p_0$.

\paragraph{Recursive marginal refinement} 
To close this gap, we introduce refinement transitions 
\begin{equation}
    K_{ \theta, t}\rbr{\dv\mid \dv'}, \qquad\dv'\sim p_{\theta, t},\qquad t=T, T-1,\ldots,1,
    \label{eq:refiner_transition}
\end{equation}
where each transition maps samples from the current draft distribution toward the target distribution. 
Applying $K_{\theta,t}$ to the current marginal $p_{\theta,t}$ induces a refined marginal
\begin{equation}
     p_{\theta, t-1}(\dv)={\textstyle \int} K_{\theta,t}\rbr{\dv\mid \dv'}p_{\theta, t}(\dv')\,d\dv'.
    \label{eq:refined_marginal}
\end{equation}
Recursive application of these transitions yields a sequence of model-induced marginals
which progressively closes the gap to the target distribution $p_0$. 
Unlike conventional diffusion, this path is not prescribed by a hand-designed forward corruption process, but instead emerges from the learned refinement dynamics themselves.

\paragraph{Refinement learning objective}
Given the current draft marginal $p_{\theta,t}$, the ideal refinement transition should make the next marginal $p_{\theta,t-1}$ closer to $p_0$. 
Plugging Eq.~\eqref{eq:refined_marginal} into Eq.~\eqref{eq:init_kl} gives
\begin{equation}
    \min_{\theta}
    \KL\!\left(
        p_0(\dv)
        \,\middle\|\,
        {\textstyle \int}
        K_{\theta,t}(\dv\mid \dv')
        p_{\theta,t}(\dv')
        \,d\dv'
    \right).
    \label{eq:ideal_refinement}
\end{equation}
This objective requires marginalizing over drafts $\dv'$, which makes direct optimization difficult. Variational learning through evidence lower bound~(ELBO) is a potential choice~\citep{alias2017variational}, 
\begin{proposition}[Joint surrogate for marginal refinement]
\label{prop:joint_surrogate}
For any refinement kernel $K_{\theta,t}(\dv\mid\dv')$ and any coupling $\gamma_t(\dv,\dv')$ with marginals $p_0(\dv)$ and $p_{\theta,t}(\dv')$, we have
\begin{equation}
\begin{aligned}
    \KL\!\left(
        p_0(\dv)
        \,\middle\|\,
        {\textstyle \int}
        K_{\theta,t}(\dv\mid \dv')
        p_{\theta,t}(\dv')
        \,d\dv'
    \right)
    &\le
    \KL\!\left(
        \gamma_t(\dv,\dv')
        \,\middle\|\,
        K_{\theta,t}(\dv\mid\dv')p_{\theta,t}(\dv')
    \right).
    \label{eq:joint_surrogate_bound}
\end{aligned}
\end{equation}
The equivalency holds if $\gamma_t\rbr{\dv, \dv'} = p_0(\dv)q(\dv'|\dv)$, where $q(\dv'|\dv)=\frac{K_{\theta,t}(\dv\mid\dv')p_{\theta,t}(\dv')}{\int K_{\theta,t}(\dv\mid\dv')p_{\theta,t}(\dv') d\dv' }$.
\end{proposition}
However, the ELBO quickly becomes unaffordable in terms of computation and memory cost w.r.t. the model size increasing, due to the requirement of the auxiliary model for posterior $q(\dv'|\dv)$ approximation.
We hereby introduce a tractable and efficient coupling, which is inherited from the diffusion learning spirit, therefore, without explicit posterior approximation; meanwhile, without predefined forward perturbation, therefore, bypassing the presumed neighborhood in discrete space. 
Intuitively, we revisit the generation procedure: given the target $\dv$, the model is gradually refining the output $\dv'$ aiming towards $\dv$, which implies ${\gamma\rbr{\dv, \dv'}} =  {p_0\rbr{\dv}}q\rbr{\dv'|\dv}$ should reflect the reverse of such a generation procedure. 
Such a dependency can be easily induced through the current learned $p_{\theta, t}\rbr{\cdot}$. Concretely, given the context $\dy$ of $\dv$, which could be the prefix or random mask version of $\dv$ in unconditional generation or the corresponding prompts in conditional generation, we can easily construct $q(\dv'|\dv) = p_{\theta, t}\rbr{\dv'|\dy}$. For the detailed efficient implementation of the sampling operation, please refer to~\appref{app:implementation}. 
With our construction, the refinement kernel is pushing further towards the ground truth $\dv$ from current intermediate output $\dv'$, therefore, we obtain our objective,
\begin{equation}
    \min_{\theta} \mathcal{L}_{t}(\theta)
    :=
    -
    \mathbb{E}_{\dv\sim p_0, \dv'\sim p_{\theta,t}}\sbr{
    \log K_{\theta,t}(\dv\mid\dv')}.
    \label{eq:conditional_refinement_loss}
\end{equation}
This joint probability facilitates learning the marginal objective in Eq.~\eqref{eq:ideal_refinement} through the tractable surrogate established in Proposition~\ref{prop:joint_surrogate}.



After learning $K_{\theta,t}$, next draft marginal $p_{\theta,t-1}$ becomes available through Eq.~\eqref{eq:refined_marginal} and we repeat the procedure to further reduce the gap between the current draft distribution and the target distribution.
In \method, the joint $\gamma_t$ is constructed from model-generated drafts $\dv'\sim \ptil_{\theta,t}$ and clean targets $\dv\sim p_0$ under the same conditioning context, rather than from a predefined forward perturbation. 
Thus, the intermediate states used for learning are induced by the sampler itself.
This recursive marginal refinement formulation has two practical benefits.

\paragraph{Neighborhood-agnostic refinement}
By removing the explicit forward corruption process, \method does not require a predefined token-level neighborhood perturbation. 
This is especially important in discrete language spaces, where there is no canonical metric for deciding which token states should be treated as nearby or easily reversible.


\paragraph{Capacity-adaptive refinement}
In our construction, each refinement transition is trained to move the current draft distribution closer to the target distribution $p_0$. 
Since $K_{\theta,t}$ is constrained by finite model capacity and finite candidate search, each refinement step closes only the portion of the remaining distributional gap that the current parameterization can represent. 
As a result, the useful number of refinement steps adapts naturally to the expressiveness of the learned refinement model.

Furthermore, we also provide formal justifications for the good properties of recursive marginal refinement in Appendix~\ref{app:proofs}, showing that the joint surrogate upper-bounds the marginal refinement objective, preserves the target marginal at the optimum, and yields monotonic improvement under ideal refinement.

\subsection{Parameterization Design for Refinement Transitions}
\label{sec:parameterization_design}
The \method framework allows flexible parameterizations of the refinement transition $K_{\theta}$. In the following, we discuss two design choices for $K_{\theta}$.

\paragraph{Self-refinement.}
Given a current draft $\dv'$, the model predicts a clean sequence through a token-factorized conditional distribution, 
\begin{equation}
  K_{\theta}^\mathrm{SR}\rbr{\dv\mid \dv'}
    =\textstyle
    \prod_{i=1}^{L}
   K_{\theta}^\mathrm{SR}\rbr{x_i\mid \dv'},
    \label{eq:self_refine_transition}
\end{equation}
where the model conditions on the full current draft and outputs per-token refinement probabilities. Plugging \eqref{eq:self_refine_transition} into Eq.~\eqref{eq:conditional_refinement_loss} gives the self-refinement objective
\begin{equation}
   \min_{\theta}  \Lcal_{\mathrm{SR}}(\theta) :=-\EE_{t,\,\dv\sim p_0,\,\dv'_t\sim p_{\theta,t}}
    \sbr{\textstyle\sum_{i=1}^{L}\log K_{\theta}^\mathrm{SR}(x_i\mid \dv'_t)}.
    \label{eq:self_refine_loss}
\end{equation}
At inference time, the model repeatedly feeds its own predicted draft back into the same refinement operator: $\dv_{k+1} = \operatorname{Decode}\rbr{K_{\theta}^\mathrm{SR}(\cdot\mid y_k)}$, where $\operatorname{Decode}(\cdot)$ can be greedy decoding, sampling, or threshold-based commitment.  This yields a stochastic fixed-point refinement process in which high-quality samples should become increasingly stable as refinement proceeds.

\paragraph{Best-of-$N$ refinement.} 
The self-refinement parameterization in Eq.~\eqref{eq:self_refine_transition} follows a single refinement trajectory from each draft. 
To expand the search space, we further introduce a stochastic best-of-$N$ refinement parameterization. 
Given a current draft $\dv'$, the self-refinement proposer $K_{\theta,\mathrm{SR}}\rbr{\dv\mid\dv'}$ first generates $N$ candidate refined sequences,
\begin{equation}
    \cbr{\dv_i}_{i=1}^N\sim K_{\theta}^\mathrm{SR}\rbr{\cdot\mid \dv'}.
\end{equation}
We then augment the proposer with a learned scorer $\scorer\rbr{\cdot}$ that evaluates and selects the top candidate via a stochastic best-of-$N$ refinement step, 
\begin{equation}
    J=\arg\max_{j\in\sbr{N}}\cbr{\scorer\rbr{\dv_j}+\xi_j},    \qquad
    \xi_j\sim\mathrm{Gumbel}(0,1), \qquad
    \dv^\star=\dv_J.
    \label{eq:gumbel_bon}
\end{equation}
By the Gumbel-Max trick~\citep{maddison2014sampling}, the selected samples follows
\begin{equation}
 \dv^\star\sim \PP\rbr{\dv^\star=\dv_\ell\mid \cbr{\dv_i}_{i=1}^N}
    =
    \textstyle\frac{
        \exp\rbr{\tau\scorer\rbr{\xb_\ell}}
    }{
        \sum_{j=1}^{N}
        \exp\rbr{\tau\scorer\rbr{\xb_j}}
    },
    \label{eq:finite_candidate_softmax}
\end{equation}
where $\tau$ is the temperature factor. 
The finite-candidate selection rule induces a valid refinement parameterization as depicted in Proposition~\ref{prop:bon_transition}.
\begin{proposition}[Induced best-of-$N$ refinement transition]
\label{prop:bon_transition}
Without loss of generality, we assume $\dv_1$ corresponds to the best candidate. This parameterizes the refinement probability as
\begin{equation}
    \begin{aligned}
        &K_{\theta, \psi}^\mathrm{BoN}\rbr{\dv^\star=\dv_1\mid \xb'}\propto\\
        &
        K_{\theta}^\mathrm{SR}\rbr{\xb_1\mid\xb'} 
        \EE_{\xb_{2:N}\sim K_{\theta}^\mathrm{SR}\rbr{\xb_i\mid\xb'} }
        \sbr{\textstyle\frac{
        \exp\rbr{\tau\scorer\rbr{\xb_1}}
    }{
        \exp\rbr{\tau\scorer\rbr{\xb_1}}+\sum_{j=2}^{N}
        \exp\rbr{\tau\scorer\rbr{\xb_j}
    }}},
     \label{eq:refine_prob_bon}
    \end{aligned}
\end{equation}
\end{proposition}
which shows that best-of-$N$ refinement selects the self-refinement proposal toward candidates preferred by the learned scorer. 
Optimizing Eq.~\eqref{eq:conditional_refinement_loss} with the transition defined in Eq.~\eqref{eq:refine_prob_bon} leads to a tractable objective for best-of-$N$ refinement
\begin{equation}
  \min_{\psi}   \Lcal_{\mathrm{BoN}}\rbr{\psi}  := 
  \EE_{t,\,\xb_1\sim p_0,\, \cbr{\dv_i}_{i=2}^N\sim K_{\theta}^\mathrm{SR}\rbr{\cdot\mid \dv'}}\sbr{\textstyle
   - \scorer\rbr{\xb_1} + \log\sum_{j=1}^N\exp\rbr{\scorer\rbr{\xb_j}}
  }.
  \label{eq:bon_loss}
\end{equation}
During this stage, the proposal distribution $K_{\theta}^\mathrm{SR}\rbr{\cdot\mid\dv'}$ is treated with stop-gradient, so that $\Lcal_{\mathrm{BoN}}$ optimizes only the scorer head parameters. 
Through Eq.~\eqref{eq:bon_loss}, the scorer learns how much more likely a candidate is under the target distribution than under the self-refinement proposal, enabling best-of-$N$ decoding to select higher-quality refinements from parallel candidates.

\subsection{Discussions and Implementation}
\label{sec:discussion}
We discuss several key properties of \method and its connections to related refinement-based generative modeling frameworks.

\paragraph{Flexible inference}
The learned refinement chain in Eq.~\eqref{eq:conditional_refinement_loss} naturally induces a sampling procedure of the same form as the backward process in Eq.~\eqref{eq:backward_joint}. 
However, forward-free refinement also enables more flexible inference beyond a fixed-horizon diffusion sampler. 
Since each refinement stage is trained to move the current marginal $p_{\theta,t}$ closer to the target distribution $p_0$, any intermediate marginal can serve as an approximate generator when computation is limited. 
Conversely, once the learned refinement dynamics approach the target distribution, the transition approximately preserves $p_0$: $p_0(\dv_0')
    \approx
    {\textstyle \int}
    K_{\theta}\rbr{\dv_0'\mid \dv_0}
    p_0(\dv_0)
    \,d\dv_0,$
which indicates that $p_0$ is approximately stationary under the learned refinement transition, allowing additional refinement iterations beyond the training horizon. 
Thus, \method provides a controllable quality-cost tradeoff at inference time, where one can stop early for faster decoding or apply more refinement steps to improve sample quality.
Detailed empirical studies are provided in Sec.~\ref{sec:exp_inference}.

\paragraph{Connections to iterative refinement methods}
The proposed formulation is related to several broader ideas in generative modeling and search. 
Consistency models~\cite{consistency-model} learn direct maps from arbitrary noisy states $\xb_t$ to clean data $\xb_0$ and enforce agreement along a prescribed diffusion trajectory. 
\method shares the $\xb_0$-prediction perspective but eliminates the prescribed trajectory. This perspective casts \method as a forward-free consistency model, where consistency is induced by the fixed-point behavior of the learned refinement operator rather than imposed along a predefined diffusion trajectory. 
The proposed best-of-N design is also related to learning-to-search~\citep{learning-to-search}, since the refinement $K_{\theta, \psi}^\mathrm{BoN}$ learns to evaluate among model-generated candidate drafts under the same dynamics used at inference time. 
Finally, the marginal recursion in Eq.~\eqref{eq:refined_marginal} is conceptually related to power iteration over distributions, where repeated application of an operator moves a distribution toward a fixed point. 
Variational power methods~\citep{power-iteration} estimate stationary distributions of Markov chains by learning correction ratios from sampled transitions. 
From this view, the learned scorer acts as a density-ratio correction between target refinements and proposal candidates, allowing best-of-N refinement to reweight sampler-induced candidates toward the target distribution.


\paragraph{Key implementation choices}
We adopt several practical designs to make recursive refinement efficient and stable. 
\textbf{(1) Iteration-sampled training.} For each sample, we start from a blank sequence and run the proposer for $t$ refinement iterations, where $t\sim\mathrm{Unif}\rbr{\cbr{1,\ldots,K}}$ and $K$ is a hyperparameter. 
Gradients are stopped through the first $t-1$ refinement steps and applied only to the final step, covering both initial draft prediction and later-stage refinement with limited training overhead. 
\textbf{(2) Blockwise semi-causal attention.} For sequence generation, we use blockwise semi-causal attention with block size $B$, where token $i$ can attend to token $j$ if their block indices satisfy $b(j)\le b(i)$, allowing bidirectional attention within each block and causal attention across previous blocks, similar to block diffusion~\citep{block_diffusion}. 
The bidirectional within-block attention enables global refinement and scoring over the current block, while the causal structure across blocks preserves efficient left-to-right block generation. 
\textbf{(3) Efficient Best-of-$N$ scoring.} For Best-of-$N$ refinement, we generate high-quality candidates using marginal-cost beam search~(Pseudocode in Appendix~\ref{app:implementation}). 
Empirically, we select candidates using $\log K_{\theta,\mathrm{SR}}\rbr{\dv\mid\dv'}+\tau\scorer\rbr{\dv}$ to combine proposal likelihood with the learned correction and mitigate prediction noise in both proposer and scorer~(justifications for this design is available in Appendix~\ref{app:proofs}). 
\textbf{(4) Shared backbone.}
For efficiency, the proposer and scorer share the same backbone parameters, with only a lightweight scorer head added for Best-of-$N$ selection. 
Further details on the attention mask, training and inference pseudocode, scorer architecture, and hyperparameter choices are provided in Appendix~\ref{app:implementation} and Appendix~\ref{app:training_details}.

